\section{Нелинейная задача встречи с конечной тягой}

\subsection{Ограничения импульсного приближения}

В предыдущих разделах управление представлялось последовательностью точечных импульсов. В этом случае действие ДУ сводится к мгновенному изменению скорости, а положение аппарата считается неизменным. Такое приближение оправдано, когда длительность каждого манёвра мала и изменением орбитальных элементов во время работы ДУ можно пренебречь.

В задачах с двигателями малой тяги это предположение может оказаться недостаточно точным. Требуемое изменение скорости в таком случае набирается за конечное время, а манёвр выполняется на некотором участке орбиты. Во время работы двигателя продолжается как пассивное движение аппарата, так и управляемое изменение орбитальных элементов. Следовательно, вклад управления уже нельзя описывать матрицей чувствительности, вычисленной в одной точке. Необходимо учитывать интегральный эффект управляющего ускорения вдоль дуги орбиты.

Подобный подход является стандартным при моделировании и оптимизации траекторий с манёврами конечной длительности \cite{morante2021survey, scarritt2012finite}. Далее рассматривается нелинейная задача встречи в приближении манёвров конечной тяги.

\subsection{Матрица чувствительности для манёвра с конечной тягой}

Рассмотрим один манёвр с конечной тягой, выполняемый на участке траектории $L \in [L_1, L_2]$, где $L_1$ и $L_2$ --- истинные долготы начала и конца манёвра. Пусть управляющее ускорение на этом участке постоянно в ОСК и равно $\mathbf{a}_{\text{ft}} \in \mathbb{R}^{3}$. Для протяжённого манёвра воздействие распределено по конечной дуге, поэтому вклад управления определяется интегрированием локальных изменений орбитальных элементов:
\begin{equation}
    \Delta \mathbf{y} = \int_{t_1}^{t_2} B(L(t)) \mathbf{a}_{\text{ft}} dt.
\end{equation}
Кроме того, во время работы двигателя изменяется большая полуось. Это приводит к изменению среднего движения $n = \sqrt{\mu / a^3}$. Поэтому изменение большой полуоси в ходе манёвра вызывает дополнительное накопление долготы. В линейном приближении
\begin{equation}
    n(a + \Delta a) \approx n(a) + \frac{dn}{da} \Delta a, \quad \quad 
    \frac{dn}{da} = -\frac{3}{2} \frac{n}{a}.
\end{equation}
Обозначим через $\mathbf{b}_a(t) = \mathbf{e}_a B(L(t)) \in \mathbb{R}^{1 \times 3}$ строку матрицы чувствительности большой полуоси к постоянному ускорению на участке манёвра, где $\mathbf{e}_a = \left( 1, \; \mathbf{0}_{1 \times 5} \right)$. Её интегральное значение от начала манёвра до момента $t$ обозначим
\begin{equation}
    \Delta \mathbf{a}(t) = \int_{t_1}^t \mathbf{b}_a(s) ds \in \mathbb{R}^{1 \times 3}.
\label{da_b_a}
\end{equation}
Фактическое приращение большой полуоси на момент времени $t$ равно $\delta a(t) = \Delta \mathbf{a}(t) \mathbf{a}_{\text{ft}}$. Тогда дополнительный вклад в среднюю долготу записывается как
\begin{equation}
    \int_{t_1}^{t_2} \frac{dn}{da} \Delta \mathbf{a}(t) \mathbf{a}_{\text{ft}} dt.
\end{equation}
Тогда в линейном приближении полное изменение орбитальных элементов на участке конечной тяги может быть записано в виде: 
\begin{equation}
    \Delta \mathbf{y} = B_{\text{ft}} (L_1, L_2) \mathbf{a}_{\text{ft}},
\end{equation}
где матрица манёвра с конечной тягой (finite trust) определяется выражением
\begin{equation}
    B_{\text{ft}} (L_1, L_2) = \int_{t_1}^{t_2} \left[ B(L(t)) + 
    \mathbf{e}_\lambda \frac{dn}{da} \Delta \mathbf{a}(t) \right] dt, \quad \quad
    \mathbf{e}_\lambda = \left( \mathbf{0}_{1 \times 5}, 1 \right)^\top.
\label{B_ft_t}
\end{equation}
Перейдём от интегрирования по времени к интегрированию по истинной долготе. Для задачи двух тел справедливо соотношение
\begin{equation}
    \frac{dt}{dL} = \frac{\eta^3}{n w(L)^2},
\end{equation}
где $\eta = \sqrt{1 - h^2 - k^2}$, $w(L) = 1 + h \sin{L} + k \cos{L}$. Тогда матрица $B_{\text{ft}}$ принимает вид
\begin{equation}
    B_{\text{ft}} (L_1, L_2) = \frac{\eta^3}{n} \left[ \int_{L_1}^{L_2} \frac{B(L)}{w(L)^2} dL + \mathbf{e}_\lambda \int_{L_1}^{L_2} \frac{dn}{da} \frac{\Delta \mathbf{a}(L)}{w(L)^2} dL \right].
\label{B_ft}
\end{equation}
Для дальнейшего удобства представим матрицу протяжённого манёвра как сумму двух слагаемых:
\begin{equation}
    B_{\text{ft}} (L_1, L_2) = B_{\text{ft}}^{(1)} (L_1, L_2) + 
    B_{\text{ft}}^{(2)} (L_1, L_2),
\end{equation}
\begin{equation}
    B_{\text{ft}}^{(1)} (L_1, L_2) = \frac{\eta^3}{n} \int_{L_1}^{L_2} \frac{B(L)}{w(L)^2} dL, \quad \quad
    B_{\text{ft}}^{(2)} (L_1, L_2) = \mathbf{e}_\lambda \frac{\eta^3}{n} \int_{L_1}^{L_2} \frac{dn}{da} \frac{\Delta \mathbf{a}(L)}{w(L)^2} dL.
\label{B_ft_1_and_B_ft_2}
\end{equation}
Первое слагаемое отвечает за прямое действие управляющего ускорения через матрицу $B(L)$. Второе слагаемое учитывает изменение средней долготы, вызванное изменением большой полуоси во время выполнения манёвра. Явные выражения для столбцов матриц $B_{\text{ft}}^{(1)} (L_1, L_2)$ и $B_{\text{ft}}^{(2)} (L_1, L_2)$ приведены в Приложении Б.

В линейном приближении вклад $j$-го манёвра с конечной тягой в конечное отклонение записывается в виде
\begin{equation*}
    \Delta \mathbf{y}_{f, j} = \Phi (t_j^c, t_f) B_{\text{ft}} (L_j^{\text{beg}}, L_j^{\text{end}}) \mathbf{a}_{\text{ft}},
\end{equation*}
где $t_j^c$ --- момент времени, соответствующий центральной точке $j$-го манёвра $L_j^c = (L_j^{\text{beg}} + L_j^{\text{end}}) / 2$, $L_j^{\text{beg}}$ и $L_j^{\text{end}}$ --- истинные долготы начала и конца манёвра.

\subsection{Связь манёвра конечной тяги с импульсным манёвром}

Покажем, что импульсный манёвр является предельным случаем манёвра конечной тяги. Пусть в момент времени $t_c$ выполняется точечный импульс скорости $\Delta \mathbf V$. В линейном приближении его влияние на орбитальные элементы имеет вид $\Delta \mathbf{y}_{\text{imp}} = B(L(t_c)) \Delta \mathbf V$. Теперь заменим этот импульс манёвром конечной тяги на малом интервале времени $[t_1, t_2]$, а ускорение $\mathbf{a}_{\text{ft}}$ на этом участке постоянно и выбрано так, чтобы полное изменение скорости совпадало с исходным импульсом:
\begin{equation}
    \Delta \mathbf V = \int_{t_1}^{t_2} \mathbf{a}_{\text{ft}} dt = 
    \mathbf{a}_{\text{ft}} \Delta t.
\end{equation}
Влияние такого манёвра на орбитальные элементы в линейном приближении: 
\begin{equation}
    \Delta \mathbf{y}_{\text{ft}} = B_{\text{ft}} (L_1, L_2) \mathbf{a}_{\text{ft}}.
\end{equation}
Обратимся к определению матрицы (\ref{B_ft_t}) и отдельно проанлизируем поведение слагаемых в предельном переходе. Если длительность манёвра стремится к нулю, то весь участок работы двигателя стягивается к точке $t_c$. Поэтому матрица $B(L(t))$ на всём интервале манёвра стремится к $B(L(t_c))$. Тогда основной вклад принимает вид 
\begin{equation}
    \int_{t_1}^{t_2} B(L(t)) \mathbf{a}_{\text{ft}} dt 
    \xrightarrow{\Delta t \rightarrow 0} 
    B(L(t_c)) \int_{t_1}^{t_2} \mathbf{a}_{\text{ft}} dt = 
    B(L(t_c)) \Delta \mathbf V.
\end{equation}
Второе слагаемое связано с изменением большой полуоси во время манёвра и с соответствующим изменением среднего движения. Используя (\ref{da_b_a}), получим
\begin{equation}
    \int_{t_1}^{t_2} \mathbf{e}_\lambda \frac{dn}{da} \Delta \mathbf{a}(t) \mathbf{a}_{\text{ft}} dt = 
    \mathbf{e}_\lambda \frac{dn}{da} \mathbf{a}_{\text{ft}} \int_{t_1}^{t_2} \left( \int_{t_1}^t \mathbf{b}_a(s) ds \right) dt.
\end{equation}
При малой длительности манёвра величину $\mathbf{b}_a(s)$ можно считать постоянной: 
\begin{equation}
    \int_{t_1}^t \mathbf{b}_a(s) ds \approx \mathbf{b}_a(t_c) (t - t_1).
\end{equation}
Тогда
\begin{equation}
    \int_{t_1}^{t_2} \mathbf{e}_\lambda \frac{dn}{da} \Delta \mathbf{a}(t) \mathbf{a}_{\text{ft}} dt \approx
    \mathbf{e}_\lambda \frac{dn}{da} \mathbf{b}_a(t_c) \mathbf{a}_{\text{ft}} \int_{t_1}^{t_2} (t - t_1) dt.
\end{equation}
Так как
\begin{equation}
    \int_{t_1}^{t_2} (t - t_1) dt = \frac{1}{2} \Delta t^2, 
\end{equation}
то
\begin{equation}
    \int_{t_1}^{t_2} \mathbf{e}_\lambda \frac{dn}{da} \Delta \mathbf{a}(t) \mathbf{a}_{\text{ft}} dt \approx
    \frac{1}{2} \mathbf{e}_\lambda \frac{dn}{da} \mathbf{b}_a(t_c) \mathbf{a}_{\text{ft}} \Delta t 
    \xrightarrow{\Delta t \rightarrow 0} 
    0.
\end{equation}
Таким образом, при стремлении длительности работы двигателя к нулю действие манёвра конечной тяги на орбитальные элементы совпадает с действием точечного импульса той же величины:
\begin{equation}
    \Delta \mathbf{y}_{\text{ft}} 
    \xrightarrow{\Delta t \rightarrow 0} 
    \Delta \mathbf{y}_{\text{imp}}.
\end{equation}
Поэтому при малой длительности манёвра импульсное решение можно использовать как хорошее начальное приближение для задачи с конечной тягой.

\subsection{Оценка ошибки перехода от импульсного манёвра к манёвру с конечной тягой}

\subsection{Управление с конечной тягой}

В приближении конечной тяги управление задаётся последовательностью участков работы двигателя. Будем считать, что на каждом участке вектор управляющего ускорения постоянен как по направлению, так и по модулю. Тогда управление можно записать в виде
\begin{equation}
    \mathbf{a}_{\text{ft}}(t) = \sum_{j=1}^{N} \Pi_j(t) a_{\text{ft}} \mathbf{e}_j,
\end{equation}
где $\Pi_j(t)$ --- индикатор $j$-го манёвра:
\begin{equation}
    \Pi_j(t)=
    \begin{cases}
        1, & t \in [t_j^{\text{beg}},\, t_j^{\text{end}}],\\
        0, & \text{иначе},
    \end{cases}
\end{equation}
а $\mathbf{e}_j \in \mathbb{R}^{3}$ задаёт направление тяги на $j$-м участке и удовлетворяет условию $\Vert \mathbf{e}_j \Vert_2 = 1$. 

Длительность $j$-го манёвра равна $\Delta t_j = t_j^{\text{end}} - t_j^{\text{beg}}$. Тогда соответствующая характеристическая скорость $\Delta V_j = a_{\text{ft}} \Delta t_j$. Поэтому при фиксированном модуле ускорения минимизация суммарной характеристической скорости эквивалентна минимизации суммарного времени работы двигателя:
\begin{equation}
    \sum_{j=1}^N \Delta V_j = a_{\text{ft}} \sum_{j=1}^N \Delta t_j.
\end{equation}

\subsection{Начальное приближение}

Построенное ранее многоимпульсное линейное решение используется для задания начального приближения в задаче с конечной тягой. 

\subsection{Постановка задачи}

\begin{equation}
\begin{aligned}
    \min_{\Delta t_j, \mathbf{e}_j, L_j^c} \quad &a_{\text{ft}} \sum_{j=1}^N \Delta t_j, \\
    \text{s.t.} \quad &\dot{\mathbf{y}} = \mathbf{f}(t, \mathbf{y}) + \sum_{j=1}^N B_{\text{ft}} (t_j^{\text{beg}}, t_j^{\text{end}}) \Pi_j(t) a_{\text{ft}} \mathbf{e}_j, \\
    &\mathbf{y}(t_i) = \mathbf{y}_i, \\
    &\mathbf{y}(t_f) = \mathbf{y}_t, \\
    &\Delta t_j \leq \Delta t_{\text{max}}, \quad j = 1, \dots, N.
\end{aligned}
\label{problem_finite_thrust}
\end{equation}
